%
\documentclass[11pt,letter]{article}
\usepackage[left,mathlines]{lineno} % line numbers
% Import the geometry package
\usepackage[top=1in, bottom=1in, left=0.92in,right=0.92in]{geometry}
\usepackage{graphicx} % Required for inserting images
%\usepackage{helvet}
\usepackage{mathtools}
\usepackage{amssymb}
%\usepackage[T1]{fontenc}



\usepackage{color}
\usepackage[table]{xcolor} 
\usepackage{tikz}          % for drawing crosses

\usepackage{amsmath}
\usepackage{amsthm} % this is for \autoref stuff
\usepackage{verbatim}
\usepackage[english]{babel}
% \usepackage{indentfirst} % suppress indent after section heading
\usepackage{xspace}
\usepackage{amsfonts}
\usepackage{mathrsfs}


\usepackage{mathtools}
\usepackage{tabularx}
\usepackage{array}
\usepackage{enumitem}
\usepackage{listings}
\usepackage{url}
\usepackage{tikz}
\usepackage[backref]{hyperref} 
\usepackage{cleveref}
\usepackage[ruled,linesnumbered]{algorithm2e}
%\usepackage{etry}
%\usepackage{authblk}
\usepackage[toc,page]{appendix}
%\geometry{a4paper} % not needed
%\usepackage{fullpage}
\usepackage{cleveref} % fix autoref problems, use \Cref always.
\usepackage{turnstile}
\usepackage[normalem]{ulem} % for strikout text using \sout{<xxx>}
\sloppy

\usepackage{appendix}
\usepackage[sanserif,full]{complexity}

% 
%   If you want to change the font the labels appear in, redefine the 
%    \showlabelfont
%    command, which by default expands to \small\ttfamily. For %    
%   example, to have
%   labels in a slanted font, and green, you could include the 
%   defnition \renewcommand{\showlabelfont}
%   {\small\slshape\color{green}}
%   in the preamble of your document (as long as you have loaded the
%   fcolorg package, of course).
%\newcommand{\showlabelfont}{\tiny\color{gray}}


% Enable watermark from this point onward




%\DeclareEmphSequence{\bfseries\itshape}


\definecolor{airforceblue}{rgb}{0.36, 0.54, 0.76}
\definecolor{darkcerulean}{rgb}{0.03, 0.27, 0.49}
\definecolor{darkelectricblue}{rgb}{0.33, 0.41, 0.77}
\hypersetup{colorlinks = true,
        linkcolor =  airforceblue,
        urlcolor = blue,
        citecolor = darkelectricblue,
        pdfstartview=FitH}

\newtheorem{theorem}{Theorem}[section]
\newtheorem*{theorem*}{Theorem}
\newtheorem{lemma}[theorem]{Lemma}
\newtheorem{claim}[theorem]{Claim}
\newtheorem*{claim*}{Claim}
\newtheorem{definition}[theorem]{Definition}
\newtheorem{proposition}[theorem]{Proposition}
\newtheorem*{proposition*}{Proposition}
\newtheorem{corollary}[theorem]{Corollary}
\newtheorem*{corollary*}{Corollary}
\newtheorem{remark}[theorem]{Remark}
\newtheorem*{remark*}{Remark}
\newtheorem{conjecture}[theorem]{Conjecture}



\newcommand{\para}{\paragraph}

%%%%%%%%%%%%%%%%%%%%%%%%%%%%%%%%%%%%%%%%
%%  Comment:  autoref names and CleverREF   
%%  As long as the lemma and definition 
%%  environments are defined 
%%  after the hyperref and cleveref packages 
%%  are loaded, using 
%%  \cref in lieu of \autoref should give 
%%  you what you want.
%%%%%%%%%%%%%%%%%%%%%%%%%%%%%%%%%%%%%%%%
 
% Original Jiaqi: 
%
%\theoremstyle{plain}
% \newtheorem{theorem}{Theorem}[section]
% \newtheorem{conjecture}{Conjecture}
% \newtheorem{lemma}[theorem]{Lemma}
% \newtheorem{proposition}[theorem]{Proposition}
% \newtheorem{corollary}[theorem]{Corollary}
% \newtheorem{maintheorem}[theorem]{Main Theorem}

%\theoremstyle{definition} % definitions should better be italicised 
 
\theoremstyle{remark} % To make proof sketch body text non italicised 
%\newtheorem*{remark}{Remark}
\newtheorem*{proofsketch}{Proof sketch}

\makeatletter
\DeclareRobustCommand\onedot{\futurelet\@let@token\@onedot}
\def\@onedot{\ifx\@let@token.\else.\null\fi\xspace}

\def\eg{\emph{e.g}\onedot} \def\Eg{\emph{E.g}\onedot}
\def\ie{\emph{i.e}\onedot} \def\Ie{\emph{I.e}\onedot}
\def\cf{\emph{c.f}\onedot} \def\Cf{\emph{C.f}\onedot}
\def\etc{\emph{etc}\onedot} \def\vs{\emph{vs}\onedot}
\def\wrt{w.r.t\onedot} \def\dof{d.o.f\onedot}
\def\etal{\emph{et al}\onedot} \def\resp{\emph{resp}\onedot}

\makeatother

\newcommand{\probP}{\text{I\kern-0.15em P}}


%complexity command
\newcommand{\PHP}{\mathrm{PHP}} %Pigeonhole Principle
\newcommand{\RankP}{\mathrm{RankP}} %Rank principle
\newcommand{\OT}{\mathrm{OT}} %odd town principle
\newcommand{\TRankP}{\mathrm{TRankP}} %Tensor Rank Principle
\newcommand{\IRankP}{\mathrm{IRankP}} % iterated rank principle
\newcommand{\IRankPE}{\mathrm{IRankPE}} %iterated rank principle with extension variables
\newfunc{\depth}{Depth}
\newfunc{\sdeg}{sdeg}
\newcommand{\sgn}{{\rm sgn}}
\newcommand{\perm}{{\rm perm}}
\newcommand{\IMM}{\ComplexityFont{IMM}}
\newfunc{\coeff}{coeff}
\newfunc{\Size}{Size}
\newcommand{\alg}{\ComplexityFont{alg}}
\renewcommand{\LIN}{\ComplexityFont{LIN}}
\newfunc{\ac}{ac}
\newcommand{\AUB}{\mathrm{AUB}} %algebrai circuit upper bound formula
\newfunc{\refute}{ref}
\newfunc{\cnf}{cnf}
\newfunc{\ecnf}{ecnf}
\newcommand{\Circuit}{%
    {
        \ComplexityFont{Circuit}
    }
}
\newcommand{\ACircuit}{%
    {
        \ComplexityFont{ACircuit}
    }
}
\newcommand{\Poly}{%
    {
        \ComplexityFont{Poly}
    }
}
\newcommand{\IPS}{%
    {
        \ComplexityFont{IPS} 
    }
}
\newcommand{\CNF}{ %
    {
        \ComplexityFont{CNF}
    }
}
\newcommand{\VAL}{ %
    {
        \ComplexityFont{VAL}
    }
}
\newcommand{\arit}{ %
    {
        \ComplexityFont{arit}
    }
}
\newcommand{\carry}{ %
    {
        \ComplexityFont{CARRY}
    }
}
\newcommand{\add}{ %
    {
        \ComplexityFont{ADD}
    }
}
\newcommand{\addition}{ %
    {
        \ComplexityFont{Addition}
    }
}
\newcommand{\modular}{ %
    {
        \ComplexityFont{Modular}
    }
}
\newcommand{\mult}{ %
    {
        \ComplexityFont{MULT}
    }
}
\newcommand{\XOR}{ %
    {
        \ComplexityFont{XOR}
    }
}
\newcommand{\SLP}{ %
    {
        \ComplexityFont{SLP}
    }
}
\newcommand{\cd}{ %
    {
        \ComplexityFont{c.d.}
    }
}



\newcommand{\PC}{ %
    {
        \ComplexityFont{PC}
    }
}
\newcommand{\PCR}{ %
    {
        \ComplexityFont{PCR}
    }
}

\newcommand{\cdIPS}{%
    {
        \ComplexityFont{cd}\text{-}\ComplexityFont{IPS}
    }
}


\newcommand{\SCNF}{%
    {
        \ComplexityFont{SCNF}
    }
}
\newcommand{\UBIT}{%
    {
        \ComplexityFont{UBIT}
    }
}

% This is changed back to newcommands instead of the dangerous "renewcommand"
\newcommand{\x}{\mathbf{x}}
\newcommand{\Z}{\mathbb{Z}}
\newcommand{\vv}{\mathbf{v}}
\newcommand{\bR}{\mathbb{R}}
\newcommand{\F}{\mathbb{F}}
\newcommand{\cM}{\mathcal{M}}
%\newcommand{\cP}{\mathcal{P}}
\newcommand{\N}{\mathcal{N}}
\newcommand{\bG}{\mathbb{G}}
%\newcommand{\cL}{\mathcal{L}}
\newcommand{\T}{\mathcal{T}}
\newcommand{\scnf}{
    {
        \ComplexityFont{scnf}
    }
}
\newcommand{\VAC}{
    {
        \ComplexityFont{VAC}
    }
}
\newcommand{\calM}{\mathcal{M}}

%% ST25 Macros for IPS proofs with better looking VDash
\newcommand{\ipspp}{\ensuremath{\hbox{\rm IPS}_{\textit{ref}}}}
\newcommand{\oldipspp}{\ensuremath{\hbox{\rm IPS}_{\textit{ref}}}}
\newcommand{\IPSnb}{\ensuremath{\IPS^{\rm alg}}}
\newcommand{\ipsppnb}{\ensuremath{\hbox{\rm IPS}^{{\rm alg}}_{\textit{ref}}}}
\newcommand{\oldipsprf}[1]{{\ensuremath{\vdash^{#1}_{\rm IPS}}}}
\newcommand{\ipsprf}[1]{{\ensuremath{\,\sststile[ss]{\mathrm{IPS}}{#1}}\,}}
\newcommand{\nipsprf}[1]{{\ensuremath{\,\sststile[ss]{\mathrm{IPS}}{#1}}\,}}
\newcommand{\ipsprfnb}[1]{{\ensuremath{\,\sststile[ss]{\mathrm{IPS^{\rm alg}}}{#1}}\,}}


% No spaces in the below macros, to avoid forcing spaces after the symbols like in the above macros. 
\newcommand{\ACZ}{\ensuremath{\ComplexityFont{AC}^0}}
\newcommand{\ACZP}{\ensuremath{\ACZ[p]}}
\newcommand{\demph}[1]{\textbf{\emph{#1}}}
\newcommand{\VNCTwo}{\ensuremath{\ComplexityFont{VNC^2}}}
\newcommand{\NCTwo}{\ensuremath{\ComplexityFont{NC^2}}}
\newcommand{\NCOne}{
\ensuremath{\ComplexityFont{NC^1}}}
        

% Commentout macro
\newcommand{\commentout}[1]{}

\let\etalchar\relax % some bug


%%%%%%%%%%%%%%%%%%%%%%%%%%%%%%%%%%%%%%%%%%%%%%%%%%%%%%%%%%%%%%%%%%%%%%
\title{\bf \ACZ$[p]$-Frege Cannot Efficiently Prove that Constant-Depth Algebraic Circuit Lower Bounds are Hard}
%%%%%%%%%%%%%%%%%%%%%%%%%%%%%%%%%%%%%%%%%%%%%%%%%%%%%%%%%%%%%%%%%%%%%%


\author{
 Jiaqi Lu\thanks{Department of Computing. This project has received funding from the European Research Council (ERC) under the European Union’s Horizon 2020 research and innovation programme (grant agreement No 101002742, \textit{EPRICOT} project).}
 \smallskip
 \\
 \small Imperial College London
 \and 
 Rahul Santhanam%
 \thanks{This project was partly funded by the EPSRC grant EP/Z534158/1, \textit{Integrated Approach to Computational Complexity: Structure, Self-Reference and Lower Bounds}.}
 \smallskip
 \\
 \small University of Oxford
 \and
 Iddo Tzameret\thanks{Department of Computing. This project has received funding from the European Research Council (ERC) under the European Union’s Horizon 2020 research and innovation programme (grant agreement No 101002742, \mbox{\textit{EPRICOT}} project). It was also supported by the Engineering and Physical Sciences Research Council (EPSRC) under grant EP/Z534158/1, \textit{Integrated Approach to Computational Complexity: Structure, Self-Reference and Lower Bounds}.. Email: \text{iddo.tzameret@gmail.com}}
  %\smallskip 
 \\ 
  \small Imperial College London
 }


% Comments ------
\newcommand{\iddo}[1]%{}
 {%
    [%
        {   
                {%
            \textcolor{green}%
                    {\footnotesize{Iddo: #1}}%
                }%
        }%
    ]%
}
\newcommand{\Iddo}{\iddo}

\newcommand{\rahul}{\Rahul}
\newcommand{\Rahul}[1]%
{%
    [%
        {   
                {%
            \textcolor{green}%
                    {\footnotesize{Rahul: #1}}%
                }%
        }%
    ]%
}
\newcommand{\jiaqi}{\Jiaqi}
\newcommand{\Jiaqi}[1]%{}
{%
    [%
        {   
                {%
            \textcolor{green}%
                    {\footnotesize{Jiaqi: #1}}%
                }%
        }%
    ]%
}
\newcommand{\new}[1]{}%{{{#1}}}

%\date{September 6, 2025}


% Change line number color
\renewcommand\linenumberfont{\color{lightgray}\tiny}

%%%%%%%%%%%%%%%%%%%%%%%%
\usepackage{fontspec}
